\chapter{Problem Statement}
\label{ch:problem_statement}

In the professional networking landscape, students and graduates of Rajiv Gandhi University of Knowledge Technologies (RGUKT) face unique structural hurdles. To understand the need for a specialized university network, we must analyze the existing communication deficits and operational challenges.

\section{Analysis of Current Communication Deficits}
Currently, the RGUKT community relies on general-purpose social networks (such as LinkedIn, Facebook, and Twitter) or private messaging channels (such as WhatsApp and Telegram) to manage alumni-student communication. These general-purpose channels exhibit major limitations:
\begin{itemize}
    \item \textbf{Fragmentation:} Unofficial messaging groups are restricted by size limits (e.g., WhatsApp group limits) and split by batches or branches. This splits the network into isolated clusters, preventing unified outreach.
    \item \textbf{High Signal-to-Noise Ratio:} General-purpose groups are often saturated with non-professional discussions, meme sharing, and administrative announcements, making it difficult for users to isolate high-value career discussions or referral posts.
    \item \textbf{Searchability Issues:} Messaging applications do not structure profile data. A student cannot run a query to find a senior who graduated in 2021 with a CSE degree and currently works as a cloud engineer at a specific target company.
\end{itemize}

\section{The Verification Gap and Trust Deficit}
A critical vulnerability in public networking websites is the absence of verification controls for institutional associations:
\begin{itemize}
    \item \textbf{Affiliation Spoofing:} Anyone can create a public profile on LinkedIn and assert that they are a student or graduate of RGUKT. This lack of validation introduces a security vulnerability, opening the door for phishing and social engineering.
    \item \textbf{Spamming Alumni:} Because public networks are open, alumni are often flooded with generic connection requests and cold messages from strangers outside their university circle. This causes network fatigue, causing alumni to ignore legitimate student inquiries.
\end{itemize}

\section{The Job Referral Bottleneck}
Getting job interviews through standard resume portal applications has become increasingly difficult due to automated Applicant Tracking System (ATS) filtering. Direct referrals from employees bypass these filters and guarantee manual review. However, the path to a referral is blocked:
\begin{itemize}
    \item \textbf{Lack of Direct Context:} Cold emailing or cold messaging alumni on LinkedIn without a shared context is often ignored.
    \item \textbf{Friction in Posting:} Alumni who want to help their juniors by referring them for active openings in their companies have no central university dashboard to post these listings. They must post on personal feeds, where reach is limited.
\end{itemize}

\section{The Proposed Solution: RGUKT Connect}
RGUKT Connect resolves these challenges by establishing a closed-loop, verified network restricted strictly to the RGUKT ecosystem. By verifying accounts through university email domains (`@rgukt.ac.in`) and university-issued student registration IDs, the platform guarantees that every user is authentic. It integrates a structured directory, direct job and referral postings, a social feed, and real-time private messaging in a single, secure space.
